% 由 work/gen_index.py 生成，勿手改
\chapter*{译注一览}
\addcontentsline{toc}{chapter}{译注一览}
\markboth{译注一览}{译注一览}
\begin{multicols}{2}\small
\noindent\textbf{1}\;\textit{含糊点拨}\;几乎同构：idele 类群的单位分支\hfill p.\,\pageref{tn:1}\\[1pt]
\noindent\textbf{2}\;\textit{含糊点拨}\;两次修改的交换性条件即余圈条件\hfill p.\,\pageref{tn:2}\\[1pt]
\noindent\textbf{3}\;\textit{易错预警}\;上标「心」记号指 t-结构的心\hfill p.\,\pageref{tn:3}\\[1pt]
\noindent\textbf{4}\;\textit{背景直觉}\;幂零奇异支撑与 AG 修正\hfill p.\,\pageref{tn:4}\\[1pt]
\noindent\textbf{5}\;\textit{含糊点拨}\;「无限级」指形式圆盘上的平凡化\hfill p.\,\pageref{tn:5}\\[1pt]
\noindent\textbf{6}\;\textit{背景直觉}\;卷积即融合积，闭包上的 IC 层对应不可约表示\hfill p.\,\pageref{tn:6}\\[1pt]
\noindent\textbf{7}\;\textit{含糊点拨}\;相对位置与 Birkhoff 分解\hfill p.\,\pageref{tn:7}\\[1pt]
\noindent\textbf{8}\;\textit{易错预警}\;移位与 Tate 扭转两个因子\hfill p.\,\pageref{tn:8}\\[1pt]
\noindent\textbf{9}\;\textit{含糊点拨}\;半和正根在基本权上的值\hfill p.\,\pageref{tn:9}\\[1pt]
\noindent\textbf{10}\;\textit{含糊点拨}\;有理平凡化与缺陷除子\hfill p.\,\pageref{tn:10}\\[1pt]
\noindent\textbf{11}\;\textit{易错预警}\;特征标群是 Pontryagin 对偶\hfill p.\,\pageref{tn:11}\\[1pt]
\noindent\textbf{12}\;\textit{含糊点拨}\;mirabolic 子群与两条轨道\hfill p.\,\pageref{tn:12}\\[1pt]
\noindent\textbf{13}\;\textit{含糊点拨}\;沿 Hecke 对应的卷积\hfill p.\,\pageref{tn:13}\\[1pt]
\noindent\textbf{14}\;\textit{含糊点拨}\;Artin–Schreier 层与洁净判别\hfill p.\,\pageref{tn:14}\\[1pt]
\noindent\textbf{15}\;\textit{背景直觉}\;标量自同构胚叠与刚性化\hfill p.\,\pageref{tn:15}\\[1pt]
\noindent\textbf{16}\;\textit{含糊点拨}\;叠的维数为负的算法\hfill p.\,\pageref{tn:16}\\[1pt]
\noindent\textbf{17}\;\textit{含糊点拨}\;ρ 扭曲与紧化整体版本\hfill p.\,\pageref{tn:17}\\[1pt]
\noindent\textbf{18}\;\textit{含糊点拨}\;归一化与反常 t-正合\hfill p.\,\pageref{tn:18}\\[1pt]
\noindent\textbf{19}\;\textit{含糊点拨}\;Deligne–Abel–Jacobi 消失\hfill p.\,\pageref{tn:19}\\[1pt]
\noindent\textbf{20}\;\textit{含糊点拨}\;两个消失界的关系\hfill p.\,\pageref{tn:20}\\[1pt]
\noindent\textbf{21}\;\textit{含糊点拨}\;「相对」Fourier–Deligne 变换\hfill p.\,\pageref{tn:21}\\[1pt]
\noindent\textbf{22}\;\textit{易错预警}\;用上叹号拉回的收缩原理\hfill p.\,\pageref{tn:22}\\[1pt]
\noindent\textbf{23}\;\textit{含糊点拨}\;基础 Whittaker 层\hfill p.\,\pageref{tn:23}\\[1pt]
\noindent\textbf{24}\;\textit{含糊点拨}\;因子化恒等式与挠商\hfill p.\,\pageref{tn:24}\\[1pt]
\noindent\textbf{25}\;\textit{含糊点拨}\;因子化恒等式即层化陈述\hfill p.\,\pageref{tn:25}\\[1pt]
\noindent\textbf{26}\;\textit{含糊点拨}\;标量胚叠与对称群的下降\hfill p.\,\pageref{tn:26}\\[1pt]
\noindent\textbf{27}\;\textit{含糊点拨}\;中间延拓的判别\hfill p.\,\pageref{tn:27}\\[1pt]
\noindent\textbf{28}\;\textit{含糊点拨}\;幂零锥商叠与仿射直线上的挠层\hfill p.\,\pageref{tn:28}\\[1pt]
\noindent\textbf{29}\;\textit{含糊点拨}\;按分层计算 Euler 示性数\hfill p.\,\pageref{tn:29}\\[1pt]
\noindent\textbf{30}\;\textit{含糊点拨}\;FGKV 比较定理的两边\hfill p.\,\pageref{tn:30}\\[1pt]
\noindent\textbf{31}\;\textit{易错预警}\;相对维数与半扭转两个因子\hfill p.\,\pageref{tn:31}\\[1pt]
\noindent\textbf{32}\;\textit{含糊点拨}\;真抛物常数项与尖点性\hfill p.\,\pageref{tn:32}\\[1pt]
\noindent\textbf{33}\;\textit{含糊点拨}\;沿满射的余极限\hfill p.\,\pageref{tn:33}\\[1pt]
\noindent\textbf{34}\;\textit{含糊点拨}\;Cauchy 分解与 Schur 函子\hfill p.\,\pageref{tn:34}\\[1pt]
\noindent\textbf{35}\;\textit{含糊点拨}\;平均函子的双侧伴随\hfill p.\,\pageref{tn:35}\\[1pt]
\noindent\textbf{36}\;\textit{背景直觉}\;几何 Ramanujan 定理\hfill p.\,\pageref{tn:36}\\[1pt]
\noindent\textbf{37}\;\textit{含糊点拨}\;半扭转的定义\hfill p.\,\pageref{tn:37}\\[1pt]
\noindent\textbf{38}\;\textit{易错预警}\;奇异支撑叠不是奇点集\hfill p.\,\pageref{tn:38}\\[1pt]
\noindent\textbf{39}\;\textit{背景直觉}\;导出 Satake 与幂零奇异支撑\hfill p.\,\pageref{tn:39}\\[1pt]
\noindent\textbf{40}\;\textit{含糊点拨}\;微局部支撑的两种定义\hfill p.\,\pageref{tn:40}\\[1pt]
\noindent\textbf{41}\;\textit{含糊点拨}\;可呈示性：余极限等于极限\hfill p.\,\pageref{tn:41}\\[1pt]
\noindent\textbf{42}\;\textit{易错预警}\;左完备化不是 Ind-完备化\hfill p.\,\pageref{tn:42}\\[1pt]
\noindent\textbf{43}\;\textit{含糊点拨}\;扭转与 Hitchin 基\hfill p.\,\pageref{tn:43}\\[1pt]
\noindent\textbf{44}\;\textit{含糊点拨}\;Betti 局部到整体的一族函子\hfill p.\,\pageref{tn:44}\\[1pt]
\noindent\textbf{45}\;\textit{含糊点拨}\;Ran 范畴是余极限\hfill p.\,\pageref{tn:45}\\[1pt]
\noindent\textbf{46}\;\textit{易错预警}\;右 t-正合不能省\hfill p.\,\pageref{tn:46}\\[1pt]
\noindent\textbf{47}\;\textit{背景直觉}\;范畴迹与可对偶性\hfill p.\,\pageref{tn:47}\\[1pt]
\noindent\textbf{48}\;\textit{含糊点拨}\;幂零奇异支撑与温和部分\hfill p.\,\pageref{tn:48}\\[1pt]
\noindent\textbf{49}\;\textit{含糊点拨}\;重构映射与双 Whittaker 范畴\hfill p.\,\pageref{tn:49}\\[1pt]
\noindent\textbf{50}\;\textit{含糊点拨}\;非正则幂零锥与反温和\hfill p.\,\pageref{tn:50}\\[1pt]
\noindent\textbf{51}\;\textit{含糊点拨}\;oper 叠是 Hitchin 挠子\hfill p.\,\pageref{tn:51}\\[1pt]
\noindent\textbf{52}\;\textit{含糊点拨}\;Kostant 切片\hfill p.\,\pageref{tn:52}\\[1pt]
\noindent\textbf{53}\;\textit{含糊点拨}\;最低 L0-权\hfill p.\,\pageref{tn:53}\\[1pt]
\noindent\textbf{54}\;\textit{易错预警}\;FLE 谱空间为何更大\hfill p.\,\pageref{tn:54}\\[1pt]
\noindent\textbf{55}\;\textit{含糊点拨}\;半无限上同调 ∞/2\hfill p.\,\pageref{tn:55}\\[1pt]
\noindent\textbf{56}\;\textit{含糊点拨}\;对偶水平配方\hfill p.\,\pageref{tn:56}\\[1pt]
\noindent\textbf{57}\;\textit{含糊点拨}\;谱侧函子对粗函子的提升\hfill p.\,\pageref{tn:57}\\[1pt]
\noindent\textbf{58}\;\textit{含糊点拨}\;插入与 Chevalley 上链\hfill p.\,\pageref{tn:58}\\[1pt]
\noindent\textbf{59}\;\textit{含糊点拨}\;式 (16.4) 的四个因子\hfill p.\,\pageref{tn:59}\\[1pt]
\end{multicols}
